\documentclass[./测控三级项目.tex]{subfiles} % 使用 subfiles 文档类，指定主文档

\begin{document}
	\subsection{题目重述}
	\begin{compactenum}
		\item 设计一个基于NE5532的电压跟随器，构建仿真电路，仿真对单峰值$50$mV,频率为$1$Hz、$150$Hz、15kHz、$1.2$MHz的正弦波和方波的跟踪效果，描述特征，分析造成差异的原因。
		\item 在实验电路板上搭建1~2路电压跟随器（板上共有4路）。
		\item 用信号发生器产生$20$Hz、$300$Hz、$15$kHz、$1.2$MHz的正弦波和方波,测试跟随器的跟踪效果。将实测结果与仿真结果进行比较，描述差异，记录正弦波与方波图案，描述实测与仿真的差异随频率增加的变化。着重体会跟随器性能随信号频率升高而衰减的特点；体会信号源仪器性能对测试工作的影响。
	\end{compactenum}
	\subsection{题目分析与实施方案}
	\subsubsection{第一问}
	第一问要求设计电压跟随器，电压跟随器是电压增益为一的运算放大器，因此在本题使用这个电路。
	使用 Multisim 软件仿真对应电路，将题2中的输出信号输入搭建的电压跟随器上，并连接示波器。观察示波器波形，分析造成差异的原因。
	\subsubsection{第二问}
	第二问依据之前设计的电压跟随器电路原理图，在PCB上进行电路焊装。
	\subsubsection{第三问}
	对焊装后的电压跟随器电路进行多方面测试，并对比与仿真波形的差异。对不同峰值和频率的正弦波、矩形波信号放大测试，能全面了解电路在实际焊装后的性能表现，包括频率响应、压摆率限制以及阶跃响应等特性，进而分析实际电路与理想仿真情况的偏差原因。
	
	\subsection{理论分析与仿真}
	\subsubsection{第一问}
	电压跟随器（也称为缓冲放大器、单位增益放大器或隔离放大器）是一种运算放大器电路，其输出电压等于输入电压，其原理图如图\ref{fig:电压跟随器电路}所示。
	\begin{figure}[h]
		\centering
		\includegraphics[width=0.45\textwidth]{电压跟随器电路}
		\caption{电压跟随器电路原理图}
		\label{fig:电压跟随器电路}
	\end{figure}\par
	由于其系统增益为1，直接将输出与负输入端短接即可，所不涉及任何其他元器件。\par
	以上理论完备，在Multisim软件中进行仿真实验。\par
	首先根据图\ref{fig:电压跟随器电路}，搭建仿真电路图，将信号发生器接入电压跟随器输入端，示波器通道A接入电压跟随器输入端，示波器通道B接入电压跟随器输出端。其完成图如图\ref{fig:电压跟随器仿真原理图}所示。
	\begin{figure}[h]
		\centering
		\includegraphics[width=0.4\textwidth]{./第二题示波器现象/电压跟随器仿真原理图}
		\caption{电压跟随器仿真原理图}
		\label{fig:电压跟随器仿真原理图}
	\end{figure}\par
	该仿真电路图测得对应的仿真正弦波形和仿真方波波形如图\ref{fig:仿真电压跟随器跟随正弦波信号}和图\ref{fig:仿真电压跟随器跟随方波信号}所示。\par
	为了更好展示相位差，正弦波除了给出电压随时间改变的$V/t$图，还给出了输出比输入的$V_B/V_A$图。\par
	根据电压跟随器的输入输出特性$V_O=V_I$，$V/t$图上的输入输出曲线应当重合，$V_B/V_A$图上的曲线应为一过原点的直线。\par
	低频时的曲线都很好的说明了这点，1Hz到15kHz的正弦、方波曲线（图\ref{fig:仿真电压跟随器跟随1Hz正弦波}、图\ref{fig:仿真电压跟随器跟随1Hz方波}、图\ref{fig:仿真电压跟随器跟随150Hz正弦波}、图\ref{fig:仿真电压跟随器跟随150Hz方波}、图\ref{fig:仿真电压跟随器跟随15kHz正弦波}、图\ref{fig:仿真电压跟随器跟随15kHz方波}）基本重合。\par
	但是图\ref{fig:仿真电压跟随器跟随1.2MHz正弦波}的$V_B/V_A$图出现了明显的迟滞现象，这表明在高频区间，运放本身的动态频率响应会对信号造成影响，造成信号的相移，需要外加电路来抑制相移的产生。
	图\ref{fig:仿真电压跟随器跟随1.2MHz方波}的$V/t$图出现了明显的不重合区域，这说明运放的增益带宽积（GBP）是有限的，高频信号会减小其压摆率造成信号跟随上的问题。
	\begin{figure}[H]
		\begin{subfigure}{0.5\linewidth}
			\centering
			\includegraphics[width=0.75\linewidth]{./第二题示波器现象/示波器1}
			\includegraphics[width=0.75\linewidth]{./第二题示波器现象/示波器9}
			\caption{仿真电压跟随器跟随1Hz正弦波信号}
			\label{fig:仿真电压跟随器跟随1Hz正弦波}
		\end{subfigure}
		\begin{subfigure}{0.5\linewidth}
			\centering
			\includegraphics[width=0.75\linewidth]{./第二题示波器现象/示波器2}
			\includegraphics[width=0.75\linewidth]{./第二题示波器现象/示波器10}
			\caption{仿真电压跟随器跟随150Hz正弦波信号}
			\label{fig:仿真电压跟随器跟随150Hz正弦波}
		\end{subfigure}
	\end{figure}
	\begin{figure}[H]\ContinuedFloat
		\begin{subfigure}{0.5\linewidth}
			\centering
			\includegraphics[width=0.75\linewidth]{./第二题示波器现象/示波器3}
			\includegraphics[width=0.75\linewidth]{./第二题示波器现象/示波器11}
			\caption{仿真电压跟随器跟随15kHz正弦波信号}
			\label{fig:仿真电压跟随器跟随15kHz正弦波}
		\end{subfigure}
		\begin{subfigure}{0.5\linewidth}
			\centering
			\includegraphics[width=0.75\linewidth]{./第二题示波器现象/示波器4}
			\includegraphics[width=0.75\linewidth]{./第二题示波器现象/示波器12}
			\caption{仿真电压跟随器跟随1.2MHz正弦波信号}
			\label{fig:仿真电压跟随器跟随1.2MHz正弦波}
		\end{subfigure}
		\caption{仿真电压跟随器跟随正弦波信号}
		\label{fig:仿真电压跟随器跟随正弦波信号}
	\end{figure}
	\begin{figure}[H]
		\begin{subfigure}{0.5\linewidth}
			\centering
			\includegraphics[width=0.75\linewidth]{./第二题示波器现象/示波器5}
			\caption{仿真电压跟随器跟随1Hz方波信号}
			\label{fig:仿真电压跟随器跟随1Hz方波}
		\end{subfigure}
		\begin{subfigure}{0.5\linewidth}
			\centering
			\includegraphics[width=0.75\linewidth]{./第二题示波器现象/示波器6}
			\caption{仿真电压跟随器跟随150Hz方波信号}
			\label{fig:仿真电压跟随器跟随150Hz方波}
		\end{subfigure}
		\begin{subfigure}{0.5\linewidth}
			\centering
			\includegraphics[width=0.75\linewidth]{./第二题示波器现象/示波器7}
			\caption{仿真电压跟随器跟随15kHz正弦波信号}
			\label{fig:仿真电压跟随器跟随15kHz方波}
		\end{subfigure}
		\begin{subfigure}{0.5\linewidth}
			\centering
			\includegraphics[width=0.75\linewidth]{./第二题示波器现象/示波器8}
			\caption{仿真电压跟随器跟随1.2MHz正弦波信号}
			\label{fig:仿真电压跟随器跟随1.2MHz方波}
		\end{subfigure}
		\caption{仿真电压跟随器跟随方波信号}
		\label{fig:仿真电压跟随器跟随方波信号}
	\end{figure}\par
	
	\subsection{实际电路测试结果与数据分析}
	\subsubsection{第二问}
	由于电压跟随器不涉及其他任何元器件，因此直接焊接芯片、电源滤波电容和信号输入输出排针即可。\par
	其焊接完成照片如图\ref{fig:电压跟随器焊接完成}所示。
	\begin{figure}[h]
		\centering
		\includegraphics[width=0.35\textwidth]{电压跟随器焊接完成}
		\caption{电压跟随器焊接完成照片}
		\label{fig:电压跟随器焊接完成}
	\end{figure}\par
	经检查分立元件位置端正、放置平整、焊点表面光滑、无凸起毛刺、无虚焊，周围无助焊剂残留；多引脚器件无短路，焊点均匀一致。符合测试要求。
	\subsubsection{第三问}
	按第一题的方案连接电源，将信号发生器信号加载在电压跟随器信号输入端，将示波器通道CH1接入电压跟随器输入端，将示波器通道CH2接入电压跟随器输出端。添加除法算子计算输出通道CH2除以输入通道CH1的增益值。\par
	将电路连接完毕后，先打开线性电源输出开关，随后根据《测控仪器电路项目作业指导书》中要求，加载对应正弦信号和方波信号。\par
	电压跟随器实验实测正弦波形和实测方波波形如图\ref{fig:实际电压跟随器跟随正弦波信号}和图\ref{fig:实际电压跟随器跟随方波信号}所示。
	\begin{figure}[H]
		\begin{subfigure}{0.5\linewidth}
			\centering
			\includegraphics[width=0.8\linewidth]{./第二题示波器现象/UpoPrtsc-001}
			\caption{实际电压跟随器跟随20Hz正弦波信号}
			\label{fig:实际电压跟随器跟随20Hz正弦波}
		\end{subfigure}
		\begin{subfigure}{0.5\linewidth}
			\centering
			\includegraphics[width=0.8\linewidth]{./第二题示波器现象/UpoPrtsc-002}
			\caption{实际电压跟随器跟随300Hz正弦波信号}
			\label{fig:实际电压跟随器跟随300Hz正弦波}
		\end{subfigure}
	\end{figure}
	\begin{figure}[H]\ContinuedFloat
		\begin{subfigure}{0.5\linewidth}
			\centering
			\includegraphics[width=0.8\linewidth]{./第二题示波器现象/UpoPrtsc-003}
			\caption{实际电压跟随器跟随15kHz正弦波信号}
			\label{fig:实际电压跟随器跟随15kHz正弦波}
		\end{subfigure}
		\begin{subfigure}{0.5\linewidth}
			\centering
			\includegraphics[width=0.8\linewidth]{./第二题示波器现象/UpoPrtsc-004}
			\caption{实际电压跟随器跟随1.2MHz正弦波信号}
			\label{fig:实际电压跟随器跟随1.2MHz正弦波}
		\end{subfigure}
		\caption{实际电压跟随器跟随正弦波信号}
		\label{fig:实际电压跟随器跟随正弦波信号}
	\end{figure}
	\begin{figure}[H]
		\begin{subfigure}{0.5\linewidth}
			\centering
			\includegraphics[width=0.8\linewidth]{./第二题示波器现象/UpoPrtsc-005}
			\caption{实际电压跟随器跟随20Hz方波信号}
			\label{fig:实际电压跟随器跟随20Hz方波}
		\end{subfigure}
		\begin{subfigure}{0.5\linewidth}
			\centering
			\includegraphics[width=0.8\linewidth]{./第二题示波器现象/UpoPrtsc-006}
			\caption{实际电压跟随器跟随300Hz方波信号}
			\label{fig:实际电压跟随器跟随300Hz方波}
		\end{subfigure}
		\begin{subfigure}{0.5\linewidth}
			\centering
			\includegraphics[width=0.8\linewidth]{./第二题示波器现象/UpoPrtsc-007}
			\caption{实际电压跟随器跟随15kHz正弦波信号}
			\label{fig:实际电压跟随器跟随15kHz方波}
		\end{subfigure}
		\begin{subfigure}{0.5\linewidth}
			\centering
			\includegraphics[width=0.8\linewidth]{./第二题示波器现象/UpoPrtsc-008}
			\caption{实际电压跟随器跟随1.2MHz正弦波信号}
			\label{fig:实际电压跟随器跟随1.2MHz方波}
		\end{subfigure}
		\caption{实际电压跟随器跟随方波信号}
		\label{fig:实际电压跟随器跟随方波信号}
	\end{figure}\par
	观察算子运算结果，可以发现，跟随器的增益系数都基本稳定在1附近，说明电路搭建正确，现象正确，测试无误。\par
	对比图\ref{fig:仿真电压跟随器跟随正弦波信号}仿真的正弦信号和\ref{fig:实际电压跟随器跟随正弦波信号}实际的正弦信号，可以发现，实际正弦信号或多或少都受到了运放本身频率相应的影响造成了相移，导致输出和输入信号的零点无法完美对齐，在示波器除法运算符的计算中变成了尖峰。同时注意到，尖峰越高说明相位差越大，横向对比图\ref{fig:实际电压跟随器跟随1.2MHz正弦波}和图\ref{fig:实际电压跟随器跟随20Hz正弦波}、图\ref{fig:实际电压跟随器跟随300Hz正弦波}、图\ref{fig:实际电压跟随器跟随15kHz正弦波}，可以发现信号频率越高，尖峰越高，相位差越大，这和仿真图像说明一致。\par
	对比图\ref{fig:仿真电压跟随器跟随方波信号}仿真的方波信号和\ref{fig:实际电压跟随器跟随方波信号}实际的方波信号，可以发现，在较低频率信号的接收上上，示波器波形会有一定的变形。但是大致和仿真没有区别。在图\ref{fig:实际电压跟随器跟随1.2MHz方波}中，可以明显的看见信号发生源在高频信号的输出上，边沿也会受到压摆率的影响，造成一定的信号变形。但是通道二展示的跟随器输出信号明显具有更长的上升时间，并有明显的超调现象发生，这说明频率的提升会影响运放压摆率。同时横向对比图\ref{fig:实际电压跟随器跟随1.2MHz方波}和图\ref{fig:实际电压跟随器跟随20Hz方波}、图\ref{fig:实际电压跟随器跟随300Hz方波}、图\ref{fig:实际电压跟随器跟随15kHz方波}，从除法算子计算的到的信号上，可以明显发现随着频率的升高，运放的动态频率特性对信号输出相位的影响，这在仿真中不易看出。
	\subsection{结论与收获}
	\textbf{结论}：在较低频率信号的接收上上，示波器波形会有一定的变形。但是大致和仿真没有区别。\par
	在高频区间，运放本身的动态频率响应会对信号造成影响，造成信号的相移，需要外加电路来抑制相移的产生。而运放的增益带宽积（GBP）是有限的，高频信号会减小其压摆率造成信号跟随上的问题。总之，跟随器性能随信号频率升高而衰减。\par
	\textbf{收获}：郭嘉媛和游寅龙秉持着精益求精的态度钻研示波器，学会了直接保存示波器图案的方法，而郑嘉宁这次尝试了自主焊接零件，加强了动手能力，提高了对小零件焊接的成功次数。
	
\end{document}